\section{Ensayo de la entrevista}
\label{sec:ensayo-entrevista}

El ensayo se desarrolló como un juego de roles para explorar el proceso de venta y facturación en un establecimiento de productos eléctricos de Quevedo, en el contexto del proyecto de plataforma de ventas con recomendaciones basadas en inteligencia artificial. La participante representó a una cajera; por tanto, sus respuestas describen un escenario simulado y no constituyen testimonios obtenidos de una trabajadora real. La entrevista estructurada permitió recorrer las actividades del puesto, identificar dificultades y registrar oportunidades de investigación posterior, siguiendo el enfoque de elicitación de requisitos \cite{pohl2025,ireb2026}.

\subsection{Condiciones del ensayo}

\begingroup
\small
\renewcommand{\arraystretch}{1.15}
\begin{longtable}{@{}p{0.27\linewidth}p{0.66\linewidth}@{}}
\caption{Condiciones confirmadas del ensayo}\label{tab:ensayo-condiciones}\\
\toprule
\textbf{Aspecto} & \textbf{Registro} \\
\midrule
\endfirsthead
\toprule
\textbf{Aspecto} & \textbf{Registro} \\
\midrule
\endhead
\bottomrule
\endlastfoot
Fecha y hora de inicio & 21 de noviembre de 2026, a las 07:50 (hora de Ecuador). \\
Medio & Videollamada mediante Microsoft Teams. \\
Entrevistador & Mendoza Palma Allan Jeremy. \\
Participante & Sánchez Centeno Roselyn Andreina, en el rol simulado de cajera de un establecimiento eléctrico. \\
Responsable de grabación & Sánchez Centeno Roselyn Andreina. \\
Modalidad de registro & Audio, según lo informado por el equipo. \\
Instrumento & Guion estructurado de 13 preguntas, incluida la autorización de grabación; 12 preguntas sobre el trabajo del cajero. \\
Desarrollo técnico & El equipo informó que el ensayo transcurrió con normalidad, sin incidencias. \\
\end{longtable}
\endgroup

La duración de diez minutos fue una previsión del guion. Según lo informado por el equipo, el audio dura 5 minutos y 10 segundos. Esta duración no se ha contrastado de forma independiente con los metadatos del archivo original.

\subsection{Consentimiento}

Antes de abordar las preguntas del proceso de venta, el entrevistador explicó el propósito académico, el carácter voluntario de la participación y la posibilidad de omitir respuestas o detener la sesión. También indicó que únicamente se registraría audio con autorización expresa. La transcripción recoge la pregunta sobre la grabación de audio y la respuesta afirmativa de la participante. Además, el equipo comunicó que se obtuvo consentimiento por escrito; ese documento no ha sido examinado en esta revisión.

De acuerdo con la información proporcionada, la sesión se grabó solamente en audio. No se incorporan al repositorio público grabaciones, transcripciones íntegras ni formularios firmados. La custodia y publicación limitada de metadatos responden al tratamiento prudente de evidencias con información personal \cite{lopdp2021}.

\subsection{Desarrollo respecto del guion}

La transcripción facilitada contiene la apertura, las trece preguntas del guion y el cierre. La pregunta 1 se destinó a solicitar autorización para grabar; las preguntas 2 a 13 abordaron las actividades de caja, el proceso de venta, la búsqueda y el registro de productos, la facturación, las demoras, los errores, los momentos de alta demanda y la orientación técnica. En el texto recibido las preguntas aparecen en el orden previsto, sin omisiones identificables. La pregunta 7 se formuló en singular (``¿Qué paso de la facturación le toma más tiempo...?'') en lugar del plural utilizado en el guion, sin modificar el tema explorado. No se identifican sondeos adicionales distintos de las preguntas estructuradas.

\begingroup
\small
\renewcommand{\arraystretch}{1.14}
\begin{longtable}{@{}p{0.14\linewidth}p{0.28\linewidth}p{0.48\linewidth}@{}}
\caption{Síntesis de las respuestas del juego de roles}\label{tab:ensayo-respuestas}\\
\toprule
\textbf{Pregunta} & \textbf{Tema} & \textbf{Información registrada} \\
\midrule
\endfirsthead
\multicolumn{3}{@{}l}{\textit{Continuación de la síntesis de respuestas}}\\
\toprule
\textbf{Pregunta} & \textbf{Tema} & \textbf{Información registrada} \\
\midrule
\endhead
\bottomrule
\endlastfoot
2 & Actividades habituales & Aseo del área de caja, registro de ventas, emisión de facturas electrónicas, procesamiento de pagos y coordinación con bodega para confirmar existencias. \\
3--4 & Atención y localización & Saludo, identificación del producto, consulta de existencias, registro, cobro y entrega de factura y garantía. Búsqueda por código o nombre y confirmación física mediante personal de bodega. \\
5--6 & Registro y facturación & Uso de códigos de barras o búsqueda manual por nombre y referencia; ingreso de datos del cliente, revisión del comprobante, emisión electrónica y entrega digital o impresa. \\
7 & Demoras operativas & Verificación manual de transferencias bancarias fuera del sistema y búsqueda de productos cuyos nombres técnicos difieren de los términos utilizados por los clientes. \\
8--9 & Errores y correcciones & Posibles errores de cantidad o selección entre productos semejantes de diferente voltaje o marca. Corrección directa antes de autorizar; después de la autorización, nota de crédito y nueva factura con intervención administrativa. \\
10 & Alta demanda & Agilización del cobro y apoyo de vendedores para revisar precios o productos antes de que el cliente llegue a caja. \\
11--12 & Consulta y orientación técnica & Consulta de fichas y manuales; cuando hace falta orientación especializada, intervención de un vendedor técnico o del encargado de bodega. \\
13 & Mejora sugerida & Búsqueda que reconozca sinónimos y nombres coloquiales, además de validación de transferencias integrada al proceso de caja. \\
\end{longtable}
\endgroup

Las respuestas permiten reconocer dos dificultades recurrentes dentro del escenario representado: la identificación de productos cuando el vocabulario del cliente difiere del catálogo y la confirmación manual de transferencias bancarias. También se describen riesgos de confusión entre productos de características técnicas similares. Estas observaciones son insumos para formular preguntas de seguimiento y candidatos a requisitos; no prueban que todas esas situaciones se produzcan en un establecimiento real ni autorizan a fijar soluciones sin validación posterior.

\subsection{Manifiesto de evidencias}

El registro público de evidencias se conserva en \texttt{evidencias/manifiesto.csv}, sin publicar el audio ni el consentimiento firmado. La fecha corresponde a la realización del ensayo. El nombre del archivo y su duración de 5 minutos y 10 segundos fueron confirmados por el equipo. La huella SHA-256 se obtuvo con PowerShell sobre el archivo local y fue comunicada para su registro; no se efectuó un cotejo independiente con el audio original.

\begingroup
\small
\renewcommand{\arraystretch}{1.12}
\begin{longtable}{@{}p{0.22\linewidth}p{0.71\linewidth}@{}}
\caption{Datos declarados para el manifiesto de evidencias}\label{tab:ensayo-manifiesto}\\
\toprule
\textbf{Campo} & \textbf{Valor} \\
\midrule
\endfirsthead
\toprule
\textbf{Campo} & \textbf{Valor} \\
\midrule
\endhead
\bottomrule
\endlastfoot
Nombre de archivo & \texttt{ENV-ENSAYO.m4a} \\
Fecha del ensayo & 21 de noviembre de 2026 \\
Duración & 00:05:10 (5 min 10 s), según lo informado. \\
SHA-256 & \texttt{E4D7B29131E209E2E74E581A44878354}\newline\texttt{E5FE31C9D8A8B2A05C76109982752BB1} \\
\end{longtable}
\endgroup

